% META: {"id": "1SPE-SECDEG-EV-B-corrige", "chapitre": "1SPE-SECOND-DEGRE", "type_objet": "corrige_evaluation", "evaluation_ref": "1SPE-SECDEG-EV-B", "version": "B", "capacites_codes": ["C1", "C2", "C3", "C4", "C5", "C6", "C7", "C8"], "sources_inspiration": [], "parametres_sympy": {"a": 3, "b": -12, "c": 9, "a2": -2, "b2": 8, "c2": -6}, "fichier_tex": "chapitres/1SPE-SECOND-DEGRE/evaluations/1SPE-SECDEG-EV-B-corrige.tex", "status": "generated"}
% BEGIN-VERIFY
% from sympy import symbols, Rational, expand, solve
% x = symbols('x')
% f = 3*x**2 - 12*x + 9
% a, b, c = 3, -12, 9
% alpha = Rational(-b, 2*a)
% assert alpha == 2
% beta = f.subs(x, alpha)
% assert beta == -3
% assert expand(3*(x-2)**2 - 3) == f
% Delta = b**2 - 4*a*c
% assert Delta == 36
% x1 = Rational(12-6, 6)
% x2 = Rational(12+6, 6)
% assert x1 == 1
% assert x2 == 3
% assert expand(3*(x-1)*(x-3)) == f
% # Somme et produit des racines lus sur les coefficients (C7).
% S = Rational(-b, a)
% P = Rational(c, a)
% assert S == 4 and P == 3
% assert x1 + x2 == S and x1 * x2 == P
% # Le signe se lit sur la forme factorisee (C8) : positif hors des racines,
% # negatif entre elles, nul en 1 et en 3.
% assert f.subs(x, 0) > 0 and f.subs(x, 4) > 0
% assert f.subs(x, 2) < 0
% assert f.subs(x, x1) == 0 and f.subs(x, x2) == 0
% g = -2*x**2 + 8*x - 6
% assert expand(-2*(x-2)**2 + 2) == g
% assert g.subs(x, 2) == 2
% # Exercice4 : vrai modèle quadratique et contrainte de frontière.
% A_enclos = x*(100-2*x)
% x_opt = 25
% assert expand(A_enclos-(-2*(x-x_opt)**2+1250)) == 0
% assert 0 < x_opt < 50
% assert A_enclos.subs(x,x_opt) == 1250
% x_limite = 20
% assert 0 < x_limite < x_opt
% aire_limitee = A_enclos.subs(x,x_limite)
% assert aire_limitee == 1200
% assert expand(aire_limitee-A_enclos-(x_limite-x)*(60-2*x)) == 0
% assert 60-2*x_limite == 20 > 0
% END-VERIFY

%---------------------------------------------------------------
% CORRIGE — Évaluation B (55 min) — Second degré
% Bareme total : 20 points
%---------------------------------------------------------------

\begin{corrige}{1SPE-SECDEG-EV-B}

%===============================================================
\section*{Exercice 1 \hfill (7 points) — C1, C3, C4}
%===============================================================

\baremeIndicatif{Q1 : 1 pt (calculer) ; Q2 : 2 pts (calculer, communiquer) ; Q3 : 2 pts (calculer) ; Q4 : 2 pts (calculer, communiquer)}

\bigskip
\textbf{Question 1 — Identification des coefficients.}

\commentaireMarge{Dans $ax^2+bx+c$ : $a$, $b$, $c$ sont les coefficients.}

Pour $f(x) = 3x^2 - 12x + 9$ : $a = 3$, $b = -12$, $c = 9$.

\medskip
\textbf{Question 2 — Forme canonique.}

\commentaireMarge{Compléter le carré dans $x^2-4x$.}

On écrit :
\[
f(x)=3(x^2-4x)+9
=3\bigl((x-2)^2-4\bigr)+9=3(x-2)^2-3.
\]
On lit $\alpha=2$ et $\beta=-3$.

\medskip
\textbf{Question 3 — Discriminant et racines.}

\commentaireMarge{$\Delta = (-12)^2 - 4 \times 3 \times 9 = 144 - 108$.}

\[
\Delta = 144 - 108 = 36 > 0. \quad \sqrt{\Delta} = 6.
\]

\[
x_1 = \frac{12 - 6}{6} = 1, \qquad x_2 = \frac{12 + 6}{6} = 3.
\]

\medskip
\textbf{Question 4 — Forme factorisee.}

\commentaireMarge{$f(x) = a(x - x_1)(x - x_2)$ avec $a = 3$, $x_1 = 1$, $x_2 = 3$.}

\[
\boxed{f(x) = 3(x-1)(x-3).}
\]

Vérification : $3(x-1)(x-3) = 3(x^2 - 4x + 3) = 3x^2 - 12x + 9 = f(x)$. \checkmark

%===============================================================
\section*{Exercice 2 \hfill (5 points) — C2}
%===============================================================

\baremeIndicatif{Q1 : 2 pts (calculer, communiquer) ; Q2 : 1 pt (raisonner) ; Q3 : 1 pt (représenter) ; Q4 : 1 pt (calculer)}

\bigskip
\textbf{Question 1 — Forme canonique de $g$.}

\commentaireMarge{On developpe la forme canonique proposee et on vérifie.}

On developpe $-2(x-2)^2 + 2$ :
\[
-2(x-2)^2 + 2 = -2(x^2 - 4x + 4) + 2 = -2x^2 + 8x - 8 + 2 = -2x^2 + 8x - 6 = g(x). \quad \checkmark
\]

Donc $g(x) = -2(x-2)^2 + 2$ est bien la forme canonique de $g$.

\medskip
\textbf{Question 2 — Sommet et extremum.}

\commentaireMarge{$a = -2 < 0$ : parabole en $\cap$, le sommet est un maximum.}

Le sommet est $S(2\,;\,2)$. Comme $a = -2 < 0$, la parabole est tournee vers le bas : $g$ admet un \textbf{maximum} de valeur $g(2) = 2$, atteint en $x = 2$.

\medskip
\textbf{Question 3 — Tableau de variations.}

\begin{center}
\begin{tabular}{|c|ccccc|}
\hline
$x$ & $-\infty$ & & $2$ & & $+\infty$ \\
\hline
$g(x)$ & $-\infty$ & $\nearrow$ & $2$ & $\searrow$ & $-\infty$ \\
\hline
\end{tabular}
\end{center}

\medskip
\textbf{Question 4 — Axe de symetrie.}

L'axe de symetrie de la parabole a pour équation $x = \alpha = 2$.

%===============================================================
\section*{Exercice 3 \hfill (4 points) — C4, C5, C7, C8}
%===============================================================

\baremeIndicatif{Q1 : 1 pt (calculer, raisonner) ; Q2 : 1 pt (calculer, représenter) ; Q3 : 1 pt (calculer) ; Q4 : 1 pt (calculer)}

\bigskip
\textbf{Question 1 — Somme et produit des racines.}

\commentaireMarge{$S = -b/a$ et $P = c/a$ : deux quotients lus sur les coefficients, sans aucun discriminant.}

Pour $f(x) = 3x^2 - 12x + 9$, on a $a = 3$, $b = -12$ et $c = 9$, donc
\[
  S = -\frac{b}{a} = \frac{12}{3} = 4
  \qquad\text{et}\qquad
  P = \frac{c}{a} = \frac{9}{3} = 3.
\]
Deux entiers de somme $4$ et de produit $3$ : ce sont $1$ et $3$. Vérification :
$1 + 3 = 4$ et $1 \times 3 = 3$. Les racines sont donc $x_1 = 1$ et $x_2 = 3$.

\medskip
\textbf{Question 2 — Tableau de signes à partir de la forme factorisée.}

\commentaireMarge{Le facteur constant $3$ est positif : il ne change aucun signe, mais il doit figurer dans le tableau.}

$f(x) = 3(x-1)(x-3)$.

\begin{center}
\begin{tabular}{|c|ccccccc|}
\hline
$x$ & $-\infty$ & & $1$ & & $3$ & & $+\infty$ \\
\hline
$x - 1$ & & $-$ & $0$ & $+$ & $+$ & $+$ & \\
$x - 3$ & & $-$ & $-$ & $-$ & $0$ & $+$ & \\
$3$ & & $+$ & $+$ & $+$ & $+$ & $+$ & \\
\hline
$f(x)$ & & $+$ & $0$ & $-$ & $0$ & $+$ & \\
\hline
\end{tabular}
\end{center}

\medskip
\textbf{Question 3 — Résolution de $f(x) \leq 0$.}

D'après le tableau de signes : $\boxed{\mathcal{S} = [1\,;\,3]}$.

\medskip
\textbf{Question 4 — Résolution de $f(x) > 0$.}

$\boxed{\mathcal{S} = \,]-\infty\,;\,1[ \cup \,]3\,;\,+\infty[}$.

%===============================================================
\section*{Exercice 4 \hfill (4 points) — C6}
%===============================================================

\baremeIndicatif{Q1 : 1 pt (modéliser) ; Q2 : 1 pt (calculer) ; Q3 : 1 pt (raisonner, calculer) ; Q4 : 1 pt (chercher, raisonner)}

\bigskip
\textbf{Question 1 — Modèle et domaine.}

\commentaireMarge{La clôture ferme deux côtés de longueur $x$ et un côté parallèle au mur.}

Le côté parallèle mesure $100-2x$ mètres. Ainsi,
\[A(x)=x(100-2x)=-2x^2+100x.\]
Les dimensions sont strictement positives si $x>0$ et $100-2x>0$,
soit $0<x<50$.

\medskip
\textbf{Question 2 — Forme canonique.}

\commentaireMarge{Compléter le carré dans $x^2-50x$.}

\[A(x)=-2(x^2-50x)
=-2\bigl((x-25)^2-625\bigr)=-2(x-25)^2+1250.\]

\medskip
\textbf{Question 3 — Optimum admissible.}

Comme $(x-25)^2\geq0$, l'aire vérifie $A(x)\leq1250$ pour tout réel $x$.
L'égalité a lieu exactement pour $x=25$, qui appartient à $]0\,;\,50[$.
L'aire maximale est donc $1250$~m$^2$ ; les deux côtés perpendiculaires
mesurent $25$~m et le côté parallèle $50$~m.

\medskip
\textbf{Question 4 — Nouvelle frontière.}

La valeur $25$ n'est plus admissible car $25>20$.
Sur $]0\,;\,20]$, on compare l'aire à sa valeur au bord :
\[A(20)-A(x)=(20-x)(60-2x)\geq0.\]
En effet, $20-x\geq0$ et $60-2x\geq20>0$ sur ce domaine.
L'égalité n'a lieu qu'en $x=20$ : le nouveau maximum est
\[A(20)=20\times(100-2\times20)=20\times60=1200\text{ m}^2.\]
Les dimensions réalisables sont $20$~m pour chaque côté perpendiculaire
et $60$~m pour le côté parallèle au mur.

\end{corrige}
